\section{Corrections to versions 1--9}\label{sec:corrections}

This appendix lists the principal claims of versions 1--9 that were withdrawn or changed in
version~10, so that a reader of an earlier version can see which statements should no longer be
relied on.  It is not guaranteed
complete: where an earlier version and this one disagree, this one is intended.  Items marked \emph{wrong} are
contradicted by a computation or an argument in this version; items marked \emph{not
re-verified} may be right but were not recomputed with the exact methods of
Section~\ref{sec:methods} and are therefore no longer claimed.

\paragraph{Withdrawn as wrong.}
\begin{enumerate}[label=C\arabic*.,leftmargin=2.6em]
\item \emph{Golden-ratio minimum 52.}  A 51-ray uncolorable subset exists
(Section~\ref{sec:catalogue}).  Version~9 also reported a 58-ray set for the completed pool of
$\{0,\pm1,\pm\sqrt[3]2,\pm\sqrt[3]4\}$, from the same unsound procedure.  We have not reproduced
a 58-ray set and do not claim one; that figure is unsupported and withdrawn, not disproved.  The
smallest set we have is 60 rays.
\item \emph{``OCUS-certified'' minima.}  The certification code was unsound
(Section~\ref{sec:methods}).  All minima in this version were recomputed with a different
procedure.  The values 31, 33, 33, 33 and 43 for the integer, Peres, $\sqrt{-2}$, Eisenstein and
$\Q(\sqrt{-7})$ raw pools survive; 52 does not.
\item \emph{``Islands are isolated; no continuous path connects one KS-supporting alphabet to
another.''}  The alphabet $\{0,\pm1,\pm x,\pm x^{*}\}$ is uncolorable along the whole circle
$|x|^2=2$, which joins the Peres, $\sqrt{-2}$ and $\Q(\sqrt{-7})$ alphabets
(Theorem~\ref{thm:circle}).  The sweep of 164 angles offered as evidence could not have
established isolation in any case.
\item \emph{The norm bound}: ``generators with $|x|^2\ge3$ that are not roots of unity produce no
KS sets.''  It fails for $x=2$, the alphabet of $\CK$; $|x|^2$ is not an invariant of the pool
($x$ and $1/x$ give the same rays); and version~9 itself reported uncolorable pools from
generators of squared modulus 4, 5 and~6.
\item \emph{The two- and three-mechanism classifications} (``modulus-2'', ``phase'',
``minimal-polynomial'', ``indirect modulus-2'').  The statement that an uncolorable pool requires
a cancellation identity is true by definition.  The list of mechanisms was extended after each
counterexample and made no prediction.  What replaces it is
Theorems~\ref{thm:twosymbol}, \ref{thm:threesymbol}, \ref{thm:conj}, \ref{thm:circle}
and~\ref{thm:6n} and Table~\ref{tab:zerosums}.
\item \emph{Completeness of the table of three-term zero-sums.}  The earlier table was compiled
by inspection, its justification miscounted the cases, it omitted two-term sums, and it described
solution sets that are circles or lines as isolated points.  Table~\ref{tab:zerosums} is the
complete enumeration.  The accompanying claim that an SMT solver had certified completeness
overstated what was run.
\item \emph{``Universal merge saturation''} and the conjecture built on it.  Identifying two
non-orthogonal rays of any uncolorable configuration gives an uncolorable configuration: a
coloring of the quotient pulls back to the original by giving the two rays the same value.  The
3{,}756 SAT calls verified this lemma; nothing about minimal KS sets follows.  (The printed
definition also described deletion of a ray, not identification.)
\item \emph{The Galois-theoretic reformulation.}  The statement that $(\Z/n\Z)^{*}$ surjects onto
$\Z/6\Z$ if and only if $6\mid n$ is false in both directions ($n=6$, $n=7$); the conductor of
$\Q(\zeta_6)$ is 3, not 6; the norm characterization fails for real quadratic generators; and the
list of ``KS-supporting'' imaginary quadratic fields contradicted the paper's own tables.  The
parallel drawn with Gleason's theorem and with the matrix trace is withdrawn with it.
\item \emph{The Hoffman-bound explanation of CSW advantages}, the ``auxiliary ray'' mechanism and
the description of the Peres pool as disconnected (it is connected).
\item \emph{``$|S_A|+|S_B|\ge|\mathcal B|$ in every case''} for bipartite perfect strategies
(the $\Q(\sqrt{-7})$ row of version~9's own table violated it), and the
claim that every Bell scenario obtainable from $\CK$ is obtainable more cheaply from the
Eisenstein set, which rested on 200 heuristic trials.
\item \emph{The tables of bipartite-strategy costs and Cabello--Severini--Winter quantities.}
The bipartite costs came from a search that was not exhaustive.  Both tables are omitted from
this version, which makes no claim about either.  Recomputed values (exact minimum costs for
nine sets; independence number, Lov\'asz number and fractional packing number for nine sets and
six pools) are in the repository as \texttt{results/bpqs\_exact\_v10.jsonl} and
\texttt{results/csw\_v10.jsonl}.
\item \emph{``The Peres and $\Z[\sqrt{-2}]$ islands are isomorphic.''}  True of the 33-ray minimum
sets, false of the pools (Computed fact~\ref{fact:pools}).
\item \emph{Counts.}  $\{0,\pm1,\pm i\}$ has 31 rays, 69 pairs, 7 triads, not 25/48/6.  The
completed pools of $\sqrt3$ and $\sqrt5$ have 133 rays, not 145.  The Eisenstein pool is not
closed under completion (165 rays after completion, not 57).  The $\Q(\sqrt{-7})$ pool has 522
orthogonal pairs in total; version~9 said its 76 auxiliary rays ``add'' 522.  The integer pool
has 6 minimum sets in one symmetry class, not 162 in six.
\item \emph{The Hermitian cross-product formula} was printed with a spurious alternating sign;
the third ray of an orthogonal pair $v,w$ is $[(v\times w)^{*}]$.
\item \emph{``All coordinates are exact, with no floating-point approximation''} was not true of
all the scripts cited.  In this version it is true of the pool counts, colorability decisions,
minima and isomorphism tests; the rigidity ranks are numerical.
\item \emph{Further items.}  The claim that enlarged Eisenstein alphabets keep minimum 33 (once
the letters $\pm2$ are admitted the pool contains $\CK$).  The use of ``contextual fraction'' for
a ratio of bases, a term with an established different meaning.  The old certification script
could print ``certified'' after finding a smaller counterexample.\item \emph{Smaller points.}  The 13-ray pool was called too small ``because the minimum known is
31'': 31 is an upper bound, and the relevant fact is the lower bound 24.  Infinitesimal rigidity
was called necessary for global rigidity; it is neither necessary nor sufficient.  The rigidity
theorem of Trandafir and Cabello concerns realizations of the $\CK$ graph and does not imply that
every 31-ray KS set is $\CK$.  The Kirchweger--Peitl--Szeider reference had the wrong venue,
pages and an author's initial.
\end{enumerate}

\paragraph{No longer claimed (not re-verified).}
The survey of imaginary quadratic fields of class number 2--4; the survey of 28 two-generator
alphabets; the quartic alphabets; mixed alphabets with two independent irrationals; the
trigonometric sweep; the icosahedral directions; the quaternionic and octonionic alphabets (for
the latter the notions of ray and orthogonality were never defined); the comparison with the
97-ray closure of Trandafir and Cabello; the $N(S)$ comparison with Cortez, Morales and Reyes;
the ``graph universality'' of $\CK$ across group-orbit constructions; the ``30{,}000 random
hypergraphs''; the identification of the coordinate field of the original 117-ray set of Kochen
and Specker; the statement that the 43-ray and golden-ratio sets are new to the literature.
